% ============================================================================
% CAPÍTULO 7 — React\index{React}
% Texto completo (rodada editorial 2)
% ============================================================================
\chapter{React: Componentes e Interfaces Interativas}
\label{cap:react}

\begin{caixaconceito}
\textbf{React} é a biblioteca de interfaces mais usada do mundo
\cite{so2025,stateofjs2024}: constrói a interface como uma \emph{função do
estado} — dado entra, UI sai. O React 19 \cite{react19} estabilizou Server
Components e Actions, transformando a biblioteca em uma ferramenta full stack.
\end{caixaconceito}

Ao final deste capítulo, você será capaz de:
\begin{objetivos}
  \item Pensar em interface como uma \emph{árvore de componentes};
  \item Criar componentes\index{componentes} com JSX\index{JSX}, props\index{props} e composição;
  \item Gerenciar estado\index{estado} com \texttt{useState} e efeitos com \texttt{useEffect};
  \item Renderizar listas, lidar com eventos e controlar formulários;
  \item Elevar estado (\emph{lifting state up}) e criar hooks\index{hooks} próprios;
  \item Aplicar as regras dos hooks e entender o papel do React Compiler;
  \item Entregar o projeto do capítulo: um quadro Kanban de tarefas.
\end{objetivos}

\section{Por que React? O modelo mental}

Antes do React, o JavaScript manipulava o DOM ``na mão'': a cada mudança de
dado, você escrevia código para encontrar o elemento certo e atualizá-lo. Com
três ou quatro estados interagindo, esse código vira um emaranhado impossível
de prever.

O React inverte o jogo: você descreve \emph{como a tela deve ser para cada
estado}, e o React cuida de levar o DOM até lá. Quando o estado muda, a função
roda de novo e a interface ``recalcula''. Isso é chamado de
\emph{programação declarativa}: você declara o resultado, não os passos.

\begin{caixaconceito}
\textbf{Declarativo vs imperativo}: imperativo diz ``crie um elemento, insira
aqui, mude a classe para lá'' (passos); declarativo diz ``se o estado é X, a
tela é assim'' (resultado). O React é declarativo — e o \emph{re-render} é o
mecanismo que mantém a tela sincronizada com o estado.
\end{caixaconceito}

\section{Componentes e JSX}

Uma página React é uma árvore de componentes — funções que recebem
\textbf{props} (entrada) e retornam \textbf{JSX} (saída). Componentes começam
com letra maiúscula, são declarativos e reagem a mudanças de estado.

\begin{lstlisting}[style=livro, language=TypeScript,
  caption={Componente com props tipadas e composição.}, label={list:react-componente}]
interface CartaoProps {
  titulo: string;
  concluida: boolean;
  onAlternar: () => void;
}

function Cartao({ titulo, concluida, onAlternar }: CartaoProps) {
  return (
    <article className={concluida ? "cartao concluido" : "cartao"}>
      <h3>{titulo}</h3>
      <input
        type="checkbox"
        checked={concluida}
        onChange={onAlternar}
        aria-label={`Marcar ${titulo} como concluída`}
      />
    </article>
  );
}

export function ListaDeTarefas() {
  const tarefas = ["Estudar React", "Praticar TypeScript"];

  return (
    <div>
      {tarefas.map((titulo) => (
        <Cartao
          key={titulo}
          titulo={titulo}
          concluida={false}
          onAlternar={() => console.log("alternou:", titulo)}
        />
      ))}
    </div>
  );
}
\end{lstlisting}

\begin{caixadica}
JSX é \emph{JavaScript com aparência de HTML}: chaves \texttt{\{\}} ``escapam''
de volta para JavaScript. A regra do React é ``pouca magia'': quase tudo o que
você vê no JSX é JavaScript de verdade (\texttt{map}, ternários,
\texttt{\&\&}).
\end{caixadica}

\begin{caixaconceito}
\textbf{Props} são a ``entrada'' do componente: dados e funções passados de
pai para filho, \emph{somente leitura}. Se um componente precisa mudar algo, o
pai passa uma função (\texttt{onAlternar}) que o filho chama. Dados sobem,
eventos descem — é o fluxo unidirecional do React.
\end{caixaconceito}

\section{Estado e eventos: \texttt{useState}}

\begin{lstlisting}[style=livro, language=TypeScript,
  caption={Estado, eventos e formulários controlados.}, label={list:react-estado}]
import { useState } from "react";

export function Contador() {
  const [contagem, setContagem] = useState(0);

  return (
    <button onClick={() => setContagem((c) => c + 1)}>
      Cliques: {contagem}
    </button>
  );
}

export function Formulario() {
  const [nome, setNome] = useState("");

  return (
    <input
      value={nome}
      onChange={(e) => setNome(e.target.value)}
      placeholder="Seu nome"
    />
  );
}
\end{lstlisting}

\begin{caixaconceito}
Estado é \emph{imutável}: você nunca altera o valor, sempre o \emph{substitui}
por um novo via \texttt{setEstado}. Por isso, use a forma funcional
\texttt{setContagem((c) => c + 1)} quando o novo valor depende do anterior —
ela é imune a atualizações em lote e a closures defasadas.
\end{caixaconceito}

\begin{caixaarmadilha}
React agrupa atualizações de estado em lote. Se você escrever
\texttt{setContagem(contagem + 1)} três vezes seguidas, \texttt{contagem}
ainda vale o valor antigo nas três chamadas — o resultado é \texttt{+1}, não
\texttt{+3}. A forma funcional \texttt{(c) => c + 1} resolve: cada chamada
recebe o valor mais recente.
\end{caixaarmadilha}

\section{Listas e a chave correta}

\begin{lstlisting}[style=livro, language=TypeScript,
  caption={Renderização de listas com keys estáveis.}, label={list:react-lista}]
interface Tarefa {
  id: number;
  titulo: string;
  concluida: boolean;
}

const tarefas: Tarefa[] = [
  { id: 1, titulo: "Estudar React", concluida: false },
  { id: 2, titulo: "Praticar TypeScript", concluida: true },
];

export function Lista() {
  return (
    <ul>
      {tarefas.map((tarefa) => (
        <li key={tarefa.id}>
          <span style={{ textDecoration: tarefa.concluida ? "line-through" : "none" }}>
            {tarefa.titulo}
          </span>
        </li>
      ))}
    </ul>
  );
}
\end{lstlisting}

\begin{caixadica}
A \texttt{key} deve ser um identificador \emph{estável e único} vindo dos
dados (\texttt{id}), nunca o índice do \texttt{map} — a menos que a lista seja
estática e nunca reordene. O React usa a \texttt{key} para saber o que mudou,
foi inserido ou removido; uma key\index{key} errada causa bugs clássicos de estado
(como o campo de um input ``seguindo'' o item errado ao reordenar).
\end{caixadica}

\section{Formulários controlados}

Um \emph{formulário controlado} tem o estado do React como fonte da verdade:
o \texttt{value} vem do estado e o \texttt{onChange} atualiza. É previsível,
fácil de validar e de integrar com APIs:

\begin{lstlisting}[style=livro, language=TypeScript,
  caption={Formulário controlado com validação simples.},
  label={list:react-formulario}]
import { useState, type FormEvent } from "react";

export function FormularioDeTarefa({ aoCriar }: { aoCriar: (titulo: string) => void }) {
  const [titulo, setTitulo] = useState("");
  const [erro, setErro] = useState("");

  function aoEnviar(evento: FormEvent) {
    evento.preventDefault();
    if (titulo.trim().length < 3) {
      setErro("O título precisa de pelo menos 3 caracteres");
      return;
    }
    aoCriar(titulo.trim());
    setTitulo("");
    setErro("");
  }

  return (
    <form onSubmit={aoEnviar}>
      <label htmlFor="titulo">Nova tarefa</label>
      <input
        id="titulo"
        value={titulo}
        onChange={(e) => setTitulo(e.target.value)}
        placeholder="Ex.: Estudar hooks"
      />
      {erro && <p role="alert">{erro}</p>}
      <button type="submit">Adicionar</button>
    </form>
  );
}
\end{lstlisting}

\begin{caixaconceito}
No React, \texttt{htmlFor} substitui \texttt{for} do HTML (porque \texttt{for}
é palavra reservada de JavaScript). Tudo o que você aprendeu sobre
acessibilidade no capítulo~\ref{cap:html} vale aqui: labels associados,
\texttt{role="alert"} para erros e botões de verdade (\texttt{<button>}).
\end{caixaconceito}

\section{Efeitos: \texttt{useEffect} e o ciclo de vida}

Componentes têm um ciclo: \emph{montar}, \emph{atualizar} e \emph{desmontar}.
O \texttt{useEffect} conecta seu código a essas fases — o lugar para busca de
dados, timers e assinaturas:

\begin{lstlisting}[style=livro, language=TypeScript,
  caption={Efeitos: busca de dados e limpeza.}, label={list:react-efeito}]
import { useEffect, useState } from "react";

interface Usuario {
  id: number;
  nome: string;
}

export function PerfilDoUsuario({ id }: { id: number }) {
  const [usuario, setUsuario] = useState<Usuario | null>(null);
  const [carregando, setCarregando] = useState(true);

  useEffect(() => {
    let ativo = true; // evita atualizar estado após desmontar

    setCarregando(true);
    fetch(`/api/usuarios/${id}`)
      .then((r) => {
        if (!r.ok) throw new Error(`HTTP ${r.status}`);
        return r.json();
      })
      .then((dados: Usuario) => {
        if (ativo) setUsuario(dados);
      })
      .catch(() => {
        if (ativo) setUsuario(null);
      })
      .finally(() => {
        if (ativo) setCarregando(false);
      });

    return () => {
      ativo = false; // limpeza ao desmontar ou quando o id muda
    };
  }, [id]);

  if (carregando) return <p>Carregando...</p>;
  if (!usuario) return <p>Usuário não encontrado</p>;
  return <p>Olá, {usuario.nome}</p>;
}
\end{lstlisting}

\begin{caixaarmadilha}
A \emph{dependência} do efeito (\texttt{[id]}) decide quando ele roda.
Esquecê-la (array vazio quando deveria ter \texttt{id}) causa dados defasados;
colocar objetos/arrays criados no render causa \emph{loop infinito} (o efeito
roda, muda estado, re-render, novo objeto, efeito roda de novo...). As duas
dores de cabeça mais famosas do React têm a mesma cura: dependências mínimas e
estáveis.
\end{caixaarmadilha}

\begin{caixadica}
No React 19, o \textbf{React Compiler} automatiza parte dessa memorização: ele
otimiza re-renders automaticamente. Mesmo assim, entenda o modelo manual
primeiro — é ele que explica o que o compilador faz.
\end{caixadica}

\section{Elevação de estado e hooks customizados}

Quando dois componentes precisam \emph{compartilhar} o mesmo estado, ele
\emph{sobe} para o ancestral comum e desce via props (lifting state up\index{lifting state up}). Quando
uma lógica de estado se repete em vários lugares, ela vira um \emph{hook
próprio}:

\begin{lstlisting}[style=livro, language=TypeScript,
  caption={Hook customizado reutilizável.}, label={list:react-hook}]
import { useCallback, useEffect, useState } from "react";

// Hook próprio: estado sincronizado com localStorage
export function useLocalStorage<T>(chave: string, inicial: T) {
  const [valor, setValor] = useState<T>(() => {
    try {
      const salvo = localStorage.getItem(chave);
      return salvo ? (JSON.parse(salvo) as T) : inicial;
    } catch {
      return inicial;
    }
  });

  useEffect(() => {
    localStorage.setItem(chave, JSON.stringify(valor));
  }, [chave, valor]);

  const definir = useCallback((novo: T | ((v: T) => T)) => {
    setValor((atual) =>
      typeof novo === "function" ? (novo as (v: T) => T)(atual) : novo,
    );
  }, []);

  return [valor, definir] as const;
}

// Uso em qualquer componente:
// const [tarefas, setTarefas] = useLocalStorage<Tarefa[]>("kanban:tarefas", []);
\end{lstlisting}

\begin{caixaconceito}
Hooks são funções que começam com \texttt{use} e \emph{encapsulam estado}.
Eles seguem duas regras de ouro: (1) só chame hooks no topo do componente ou
de outro hook — nunca dentro de \texttt{if}, loops ou funções aninhadas; (2)
só chame hooks de componentes React ou de outros hooks. O React depende da
\emph{ordem constante} das chamadas para associar estado a cada hook.
\end{caixaconceito}

\section{Composição: o padrão de design do React}

Composição supera herança: em vez de ``herdar de um componente base'', você
\emph{compõe} componentes dentro de outros, passando o que quiser via props —
inclusive outros componentes (\texttt{children}):

\begin{lstlisting}[style=livro, language=TypeScript,
  caption={Composição com children.}, label={list:react-composicao}]
import type { ReactNode } from "react";

interface CaixaProps {
  titulo: string;
  acoes?: ReactNode;
  children: ReactNode;
}

export function Caixa({ titulo, acoes, children }: CaixaProps) {
  return (
    <section className="caixa">
      <header>
        <h2>{titulo}</h2>
        {acoes}
      </header>
      <div>{children}</div>
    </section>
  );
}

// Uso: qualquer conteúdo entra como children
<Caixa titulo="Tarefas" acoes={<button>+</button>}>
  <ul>{/* itens */}</ul>
</Caixa>
\end{lstlisting}

% ============================================================================
\section{Projeto do capítulo: quadro Kanban de tarefas}
\label{sec:projeto7}

\begin{caixaprojeto}
\textbf{Objetivo:} construir um quadro Kanban com três colunas (A fazer, Em
andamento, Concluído) usando React + TypeScript, sem backend.

\textbf{Critérios de aceite:}
\begin{itemize}[leftmargin=*, itemsep=0.2em]
  \item Criar, mover e excluir tarefas com estado tipado;
  \item Colunas com contagem de itens e destaque visual;
  \item Formulário controlado para nova tarefa com validação mínima;
  \item Persistência em \texttt{localStorage} via hook próprio
  (\texttt{useLocalStorage});
  \item Componentes pequenos e reutilizáveis (\texttt{Tarefa}, \texttt{Coluna}, \texttt{Quadro});
  \item Bootstrap com Vite: \texttt{npm create vite@latest kanban -- --template react-ts};
  \item Acessível: botões reais, labels e \texttt{role="alert"} nos erros.
\end{itemize}
\end{caixaprojeto}

\begin{caixadica}
O estado do Kanban é um ótimo exercício de modelagem: uma lista de tarefas com
\texttt{id}, \texttt{titulo}, \texttt{coluna} e \texttt{ordem}. Mover uma
tarefa é apenas \emph{substituir} o estado com a coluna atualizada (lembre-se
da imutabilidade!). Neste capítulo o estado é local ao navegador; no
capítulo~\ref{cap:fullstack}, o mesmo conceito vira um produto real com banco
de dados e autenticação.
\end{caixadica}

\section{Exercícios de fixação}

\begin{exercicios}
  \item O que são props e qual a diferença entre props e estado?
  \item Por que o estado deve ser atualizado de forma imutável?
  \item Escreva um componente \texttt{Lista} que recebe \texttt{itens: string[]} e renderiza uma \texttt{ul} com \texttt{key} correta.
  \item Qual a função do array de dependências do \texttt{useEffect}? Dê um exemplo de loop infinito.
  \item Explique \emph{lifting state up} com um exemplo de dois componentes que compartilham estado.
  \item Quais são as duas regras dos hooks e por que elas existem?
\end{exercicios}

\section{Desafios}

\begin{desafios}
  \item \textbf{Contador com histórico}: um contador que guarda o histórico de cliques em um array e permite desfazer (undo).
  \item \textbf{Busca com debounce}: campo de busca que filtra uma lista com debounce de 300ms (reutilize a ideia do capítulo~\ref{cap:js}).
  \item \textbf{Tema claro/escuro}: hook \texttt{useTema} que alterna classes CSS e persiste no \texttt{localStorage}.
  \item \textbf{Drag and drop}: permita arrastar tarefas entre colunas (HTML5 drag \& drop) preservando a ordem.
\end{desafios}

\section{Exercícios de nível sênior}

\begin{seniores}
  \item \textbf{Reconciliação}: explique o que acontece no DOM quando o estado muda e por que a \texttt{key} correta evita bugs de estado (input que perde foco, lista reordenada).
  \item \textbf{Memorização consciente}: use \texttt{React.memo}, \texttt{useMemo} e \texttt{useCallback} em uma lista grande, meça com o Profiler do React DevTools e explique quando cada um ajuda (e quando atrapalha).
  \item \textbf{Padrões avançados}: implemente \emph{compound components} em um componente \texttt{Tabs} (ex.: \texttt{Tabs.List}, \texttt{Tabs.Panel}) e explique o trade-off vs.\ props simples.
  \item \textbf{Suspense e Actions}: com o React 19, use \texttt{useActionState} e \texttt{useOptimistic} para um formulário que atualiza a UI antes da confirmação do servidor (com mock).
\end{seniores}

\section{Armadilhas comuns}

\begin{itemize}[leftmargin=*, itemsep=0.4em]
  \item Mutar estado diretamente (\texttt{arr.push}) em vez de criar um novo array;
  \item \texttt{useEffect} com dependências instáveis (loops infinitos);
  \item Esquecer \texttt{key} ou usar o índice em listas dinâmicas;
  \item Chamar hooks condicionalmente ou dentro de loops (regra dos hooks);
  \item Componentes gigantes de 500 linhas — quebre em componentes pequenos e reutilizáveis;
  \item Esquecer \texttt{evento.preventDefault()} no submit (a página recarrega).
\end{itemize}

\section{Resumo do capítulo}

\begin{itemize}[leftmargin=*, itemsep=0.3em]
  \item UI = função do estado; componentes recebem props e retornam JSX;
  \item \texttt{useState} para estado, \texttt{useEffect} para efeitos;
  \item Estado é imutável; atualizações em lote exigem a forma funcional;
  \item \texttt{key} estável em listas; formulários controlados para previsibilidade;
  \item Estado compartilhado sobe para o ancestral (lifting state up);
  \item Hooks customizados encapsulam lógica reutilizável (ex.: \texttt{useLocalStorage});
  \item Projeto entregue: Kanban funcional e persistente.
\end{itemize}

\section{Referências oficiais para aprofundar}

\begin{itemize}[leftmargin=*, itemsep=0.3em]
  \item Documentação oficial do React: \url{https://react.dev}
  \item React 19 (anúncio oficial): \url{https://react.dev/blog/2024/12/05/react-19}
  \item Curso oficial \emph{Learn React}: \url{https://react.dev/learn}
  \item Regras dos Hooks (ESLint): \url{https://react.dev/learn/rules-of-hooks}
  \item React DevTools: \url{https://react.dev/learn/react-developer-tools}
\end{itemize}
